\exnumber{02}
\extitle{UI II - Ex 3 - Bibliothèque de design system}
\exfiles{ex02/fichier Figma, nomme Name\_FirstName\_DeliverableName\_Date\_VersionNumber}
\exauthorize{Ceux nommés dans l'exercice}
\makeheaderfiles

\textbf{Instructions :}

Copiez 10 éléments du design system easyJet, incluant

\begin{itemize}
\item atoms
\item molecules
\item organisms
\end{itemize}

\textbf{Livrable :}

\begin{itemize}
\item Un fichier Figma comprenant 10 éléments du design system easyJet
\item Déposez le fichier Figma sur Github, nommé ainsi : « Name\_FirstName\_DeliverableName\_Date\_VersionNumber ». Exemple : « Doe\_John\_MidFiPrototype\_05242024\_V1 ».
\item Ajoutez un titre dans le fichier.
\item Soignez l'aspect global du fichier. Il doit être clair, simple et lisible. Vous pouvez regarder canva.com, sans surcharger le dessin de détails.
\end{itemize}

\textbf{Pistes :}

\begin{itemize}
\item Vous pouvez utiliser Avenir Book à la place de la police d'easyJet
\item Explorer Noun Project pour trouver des icônes
\end{itemize}

\textbf{Ressources :}

\begin{itemize}
\item Méthodologie Atomic Design \url{https://atomicdesign.bradfrost.com/chapter-2/}
\end{itemize}

